\chapterpage{controls}{Set a view and place a marker}
\begin{enumerate}
\item Open \ui{P-rho Diagram} and set pressure to Mbar in \ui{Plot Settings}. For shot 29400, Hicks Table I gives an inferred quartz density of 6.03 g/cm³ and a pressure of 4.79 Mbar.
\item To see the experimental region clearly, set the density axis to 5.4–6.8 g/cm³ and the pressure axis to 3.5–6.3 Mbar. Click \ui{Apply}.
\item In the manual scatter inputs enter 6.030 and 4.790 and click \ui{Add}. This puts the paper's inferred state on the plot for comparison with the Hamel curve.
\end{enumerate}
\pic{settings}{The marker places the value reported by Hicks on the density–pressure diagram. Both reference fits are shown in the same units.}{88mm}
\subhead{Which type of mark should you save?}
For a value you read by eye, keep the tab, quantities and units with the manual point. For a computed crossing, select its ID in the scatter controls. IM Tool can use the associated state values on other diagrams when those values are available. Save the workflow to keep the marker list, and export a figure when a colleague only needs the visual result.
